\documentclass[border=12pt]{standalone}
\usepackage[american]{circuitikz}
\begin{document}
\begin{circuitikz}[font=\sffamily, line width=0.8pt]
\ctikzset{diodes/scale=0.65}
% Input voltage and four ideal bridge switches.
\draw (-1,4.2) to[V] (-1,0);
\draw (-1,4.2) -- (3,4.2);
\draw (-1,0) -- (3,0);
\draw (1,4.2) to[nos] (1,2.1) to[nos] (1,0);
\draw (3,4.2) to[nos] (3,2.1) to[nos] (3,0);
\node[circ] at (1,2.1) {}; \node[circ] at (3,2.1) {};
% A drives the series resonant capacitor and inductor. The crossing at x=3 is not a junction.
\draw (1,2.1) -- (1.5,2.1) -- (1.5,2.5) -- (2.8,2.5) arc[start angle=180,end angle=0,radius=0.2] -- (4,2.5)
 to[C] (5.5,2.5) to[L] (7,2.5) -- (9,2.5);
\draw (3,2.1) -- (3.7,2.1) -- (3.7,0.55) -- (9,0.55);
\draw (7.4,2.5) to[L,*-*] (7.4,0.55);
% Ideal transformer with magnetizing inductance explicitly modeled in parallel.
\draw (9,2.5) to[L] (9,0.55);
\draw (10,2.5) to[L,mirror] (10,0.55);
\draw (9.43,0.9) -- (9.43,2.5) (9.57,0.9) -- (9.57,2.5);
\node[circ,scale=0.6] at (8.85,2.25) {}; \node[circ,scale=0.6] at (10.15,2.25) {};
% Secondary feeds the AC midpoints of the diode bridge.
\draw (10,2.5) -- (10.8,2.5) -- (10.8,1.7) -- (12,1.7);
\draw (10,0.55) -- (10,-0.5) -- (14.8,-0.5) -- (14.8,1.7) -- (14,1.7);
\draw (12,0) to[D] (12,1.7) to[D] (12,3.8);
\draw (14,0) to[D] (14,1.7) to[D] (14,3.8);
\draw (12,3.8) -- (18,3.8) (12,0) -- (18,0);
\node[circ] at (12,1.7) {}; \node[circ] at (14,1.7) {};
\draw (16,3.8) to[C,*-*] (16,0);
\draw (18,3.8) to[R] (18,0);
\end{circuitikz}
\end{document}
